% No preamble. Prefix pgh-. Finite compatible regulators throughout.
\section{H\"older regularity through an ordinary prefix}
\label{pgh-section}

Fix the primitive odd ratio $\ell$ and $0<\alpha<1$. Write
$N=MR=\ell^{k+s}$, $\eta=N^{-1}$ and $h=\ell^k\eta=R^{-1}$.
Physical next-coarse spacing is one. Define the physical prefix map
\begin{equation}
 P_h(A)=\frac1{ih}\log M_\ell^k(e^{i\eta A}).
 \label{pgh-prefix}
\end{equation}
The established ordinary-trajectory theorem gives a common complex
supremum ball on which $P_h$ is holomorphic, vanishes at zero and
$\|P_h(A)\|_\infty\le K\|A\|_\infty$, uniformly in $M,R$ and
compatible periodic volume. Its dependency radius is at most
$c_\ell h$ in physical coordinates. Nested centered partitions give
\begin{equation}
 P_h(T_zA)=T_zP_h(A),\qquad z\in h\mathbb Z^d.
 \label{pgh-covariance}
\end{equation}
Only these aligned translations are asserted. Use the complete norm
\[
 \|A\|_{\mathcal N_{\alpha,\eta}}
 =\|A\|_\infty+\|\nabla_\eta A\|_\infty
                         +[\nabla_\eta A]_\alpha,
\]
with physical distances and equivalent fixed-collar versions on
periodic lifts. Constants below depend on the primitive ratio,
$\alpha$ and the group, not on $M,R$ or auxiliary grouping $J$.

\begin{lemma}[Uniform complete-field prefix bound]
\label{pgh-prefix-lemma}
On a smaller fixed ball,
\begin{equation}
 \|P_h(A)\|_{\mathcal N_{\alpha,h}}
 \le C\bigl(\|A\|_{\mathcal N_{\alpha,\eta}}
                   +\|A\|_{\mathcal N_{\alpha,\eta}}^2\bigr).
 \label{pgh-prefix-bound}
\end{equation}
The map is holomorphic into the complete output space with uniform
Banach Cauchy derivative bounds.
\end{lemma}
\begin{proof}
Uniform supremum holomorphy bounds $DP_h,D^2P_h$ on a smaller
supremum ball by constants $K_1,K_2$. For $T_h=T_{he_\mu}$,
$\nabla_hA=(T_hA-A)/h$ is an average of fine gradients; its
supremum and H\"older seminorm are bounded by those of
$\nabla_\eta A$. Equation~\eqref{pgh-covariance} gives
\begin{equation}
 \nabla_hP_h(A)=\int_0^1
 DP_h\bigl(A+t(T_hA-A)\bigr)\nabla_hA\,dt.
 \label{pgh-gradient-identity}
\end{equation}
If $\|A\|_{\mathcal N_{\alpha,\eta}}\le b$, this bounds the output
gradient by $K_1b$. Translate by $z\in h\mathbb Z^d$ and subtract.
The change in direction costs $CK_1b|z|^\alpha$. The derivative
operator changes by at most
$K_2\|T_z((1-t)A+tT_hA)-((1-t)A+tT_hA)\|_\infty$.
After multiplication by $\|\nabla_hA\|_\infty$, this costs
$CK_2b^2|z|$, hence $CK_2b^2|z|^\alpha$ when $|z|\le1$.
For larger shifts the gradient supremum suffices. This proves
\eqref{pgh-prefix-bound}, including periodic lifts and short tori
with incidences counted with multiplicity. All interpolation
segments lie in the common complex ball. At each finite regulator
the map is holomorphic between finite-dimensional normed spaces;
the displayed uniform bound on the complete-field ball and
Banach Cauchy give uniform complete-output derivative bounds.
No ordinary-average physical adjoint is used.
\end{proof}

\subsection{Supported pre-gauge fields and their differences}
\label{pgh-support-section}

Let the established ordinary representative be
$A_{\mathrm o,N}(w)=T_N(A_{\mathrm c,N}(w))$, with
$\|A_{\mathrm o,N}(w)\|_{\mathcal N_{\alpha,\eta}}
\le C_{\mathrm o}\|w\|_\infty$. Define
\begin{equation}
 a_p(w)=hP_h(A_{\mathrm o,N}(w)),\qquad
 e^{ia_p}=M_\ell^k(e^{i\eta A_{\mathrm o,N}(w)}).
 \label{pgh-pregauge}
\end{equation}
This is the selected representative before the rooted current
gauge. It is not a regularity assertion for every gauge-equivalent
representative. Lemma~\ref{pgh-prefix-lemma} implies
\begin{equation}
 \|a_p(w)\|_{\mathcal N_{\alpha,h}}
                  \le Ch\|w\|_\infty.
 \label{pgh-angle-norm}
\end{equation}
In current-unit coordinates, let $\Delta_\mu$ be the unscaled
difference between adjacent $h$-lattice sites. Thus, on $\|w\|<b$,
\begin{equation}
 \|a_p\|_\infty\le Cbh,\qquad
 \|\Delta_\mu a_p\|_\infty\le Cbh^2,\qquad
 \|\Delta_\nu\Delta_\mu a_p\|_\infty
                         \le Cbh^{2+\alpha}.
 \label{pgh-scale-bounds}
\end{equation}
More generally, for an integer current-lattice shift $z$,
\begin{equation}
 |\Delta_\mu a_p(x+z)-\Delta_\mu a_p(x)|
                  \le Cbh^{2+\alpha}|z|^\alpha.
 \label{pgh-shift-bound}
\end{equation}
For physical shifts longer than one, the gradient supremum
proves the same bound. Periodic distances use compatible lifts.
The second-difference estimate uses first-gradient H\"older
regularity; it does not assert a uniform classical $C^2$ bound.

Let support labels refer to final physical unit cells and let
$\tau$ be graph Steiner length. Fix an admissible spatial exponent
$\kappa>0$ from the classical support theorem and any prescribed
finite cardinality reserve $a\ge0$. For an output root $i$, use
one complete $h$-field fiber $E_{h,i}$ with the local
$\mathcal N_{\alpha,h}$ norm on a fixed enlarged collar of $i$,
including all required difference endpoints and H\"older quotients.

\begin{proposition}[Complete-output strict supports]
\label{pgh-support-proposition}
There are $r_a,C_a>0$, independent of $M,R$ and volume, and strict
holomorphic families
\begin{equation}
 a_p|_i=\sum_{X\ne\varnothing}a_{p;i,X}(w_X),\qquad
 a_{p;i,X}(0)=0,
 \label{pgh-family}
\end{equation}
such that, for $0<r\le r_a$,
\begin{equation}
 \sup_i\sum_X e^{\kappa\tau(\{i\}\cup X)+a|X|}
       \|a_{p;i,X}\|_{E_{h,i};r}\le C_a h r.
 \label{pgh-family-bound}
\end{equation}
The linear supported part has bound $C_ahr$ and the nonlinear
remainder has bound $C_ahr^2$, after a fixed radius reduction.
\end{proposition}
\begin{proof}
Use the classical support theorem's smooth partition and input
Banach spaces of complete supported fine fields. Reconstruction
has a uniform norm $C_{\mathrm{rec}}$; no sharply restricted
H\"older field is extended. The map $P_h\circ T_N$ is uniformly
holomorphic into $E_{h,i}$ by Lemma~\ref{pgh-prefix-lemma} and
the ordinary-transfer theorem. Its row reads at most a fixed
number $N_0$ of input physical cells in a fixed radius $r_0$.
This follows from the $O(h)$ prefix range and the fixed physical
range of $T_N$.

On a common complete-field radius $\rho_1$, Cauchy and
polarization, followed by summing only the input cell slots,
give the mixed degree-$n$ majorant
\begin{equation}
 v_n\le K
 \left(\frac{eC_{\mathrm{rec}}N_0e^{\kappa r_0}}{\rho_1}\right)^n
       =Kc^n,\qquad n\ge1.
 \label{pgh-vertex-bound}
\end{equation}
The output collar can contain $O(R^d)$ bonds, but its complete
norm is one Banach output norm, taken before the support sum.
No output coordinate count occurs in \eqref{pgh-vertex-bound}.

Substitute the corrected minimizer's complete-field support
family, whose norm with the prescribed reserve is at most
$C'_ar$. Rooted Steiner trees glue at the input slot roots,
and $e^{a|\cup_qX_q|}\le\prod_qe^{a|X_q|}$. Hence the resulting
series is bounded by $\sum_{n\ge1}K(cC'_ar)^n$, at most
$2KcC'_ar$ when $cC'_ar\le1/2$. Multiplication by $h$ proves
\eqref{pgh-family-bound}. Use one global Taylor/partition assignment
and group by the union of input variable supports before restricting
complete output coefficients to collars. Their restrictions on
overlapping collars agree; arbitrary separate decompositions at
different output roots are not used. Absolute convergence in the complete
local output norm identifies the evaluated series with
\eqref{pgh-pregauge}. Separating its linear terms gives the
quadratic remainder estimate. Applying the bounded difference
maps on the enlarged collars yields the weighted family versions
of \eqref{pgh-scale-bounds}--\eqref{pgh-shift-bound}.
\end{proof}
These are finite-regulator coefficients and can retain dependence
on the fixed regulator and baseline. Volume consistency is not
asserted. A larger prescribed reserve may shrink $r_a$.

\subsection{An oriented discrete Taylor identity}
\label{pgh-taylor-section}

For a current-unit scalar component set
$D_\nu f(v)=f(v+e_\nu)-f(v)$. Choose a fixed shortest coordinate
path $\Gamma$ from $x$ to $y$. A traversed edge has positive
direction $\nu(e)$, lower positive-coordinate endpoint $v(e)$,
and traversal sign $s(e)$. Telescoping gives
\begin{align}
 f(y)-f(x)&=\sum_\nu(y_\nu-x_\nu)D_\nu f(x)
                         +\operatorname{Rem}_{xy}(f),\notag\\
 \operatorname{Rem}_{xy}(f)
 &=\sum_{e\in\Gamma}s(e)
       [D_{\nu(e)}f(v(e))-D_{\nu(e)}f(x)].
 \label{pgh-taylor}
\end{align}
A negative traversal uses the positive difference at its lower
endpoint, not at its traversal initial point. If $n=|x-y|_1$,
\eqref{pgh-scale-bounds} and \eqref{pgh-shift-bound} imply
\begin{equation}
 |a_p(y)-a_p(x)|\le Cbh^2n,\qquad
 |\operatorname{Rem}_{xy}(a_p)|
                      \le Cbh^{2+\alpha}n^{1+\alpha}.
 \label{pgh-taylor-bound}
\end{equation}
One also has the elementary bound
$|\operatorname{Rem}_{xy}(f)|\le n^2\sup|D^2f|$.
For a winding periodic path use its lifted displacement;
no global affine function on a torus is assumed. Coherent complete
coefficients justify telescoping across collars. If a subsequent
estimate keeps a fixed starting-root spatial weight along a long
path, the exponential comparison of root weights along that path
must also be budgeted; it is not omitted by the local norm bound.

These independently derived conventions supply the scale
pattern used in CMP109, pp.~285--288,
(4.17)--(4.18) and (4.30). They do not assert the subsequent
nonlinear Ward cancellation or its local carrier assignment.

\paragraph{Constant order and scope.}
The constants above are fixed at the primitive ratio before
$R$ and $J$. Quantum coordinate-kernel estimates obtained from
activities on $J$-cells can instead contain factors depending
on their dimension and on conversion of cell distance to
current distance. Finiteness of those constants does not make
them independent of $J$. When $J$ is chosen after $R$, an
expression $C(J)R^{-\alpha}$ is therefore not yet a
noncircular contraction estimate. This regularity theorem
does not classify the low-degree kernels, supply a bounded
gauge-compatible extraction, determine a beta-function sign,
or close the quantum induction.
